% 功耗约束
\section{功耗约束：从"总功耗别超标"到"每对相邻裸片的温差都要控制"}

第三类约束来自热物理：链路越多、lane越多、功耗越大。散不出去的热量会毁掉芯片。

\subsection{最粗糙的模型：全局功率密度}

每条链路$e$的物理层功耗为$P_{\text{lane}} \cdot L_e \cdot B/R_e$。
总功耗不能超过冷却方案的散热能力：

\begin{equation}\label{eq:thermal-l0}
    \sum_{e \in \mathcal{E}} L_e \cdot \frac{P_{\text{lane}} \cdot B}{R_e} \le A_{\text{total}} \cdot q_{\max}
\end{equation}

$q_{\max}$取决于冷却方案——一条不等式，O(1)计算。

但全局功率密度忽略了热的空间分布。
晶圆上各区域的实际温度取决于局部功耗密度和横向热传导——
两个相邻裸片，一个全速运转（50W），一个空闲（5W），即使晶圆平均功耗很低，
"热点"裸片的温度也可能超限，且两者之间的温差可能导致$\mu$bump机械失效。

\subsection{更精细的模型：稳态热网络、温差与翘曲}

更精细的分析需要求解晶圆上的稳态热传导方程：
$\nabla \cdot (k \nabla T) + \dot{q} = 0$。

\textbf{关键数学事实}：这是一个\textbf{线性偏微分方程}——温度$T$对热源$\dot{q}$的依赖天然是线性的。
离散化为热网络后：$\mathbf{T} = \mathbf{G}^{-1}(\mathbf{P}_0 + \mathbf{C} \cdot \mathbf{L} + \mathbf{b})$，
其中$\mathbf{G}$为热导矩阵，$\mathbf{P} = \mathbf{P}_0 + \mathbf{C} \cdot \mathbf{L}$为各节点功耗（$L_e$的线性函数）。

由此得到：
\begin{itemize}
    \item \textbf{温度上限约束}：$T_v \le T_{\max}$对所有裸片$v$成立——
    由于$\mathbf{T}$是$\mathbf{L}$的线性函数，这是$\mathbf{L}$上的一组线性不等式。
    \item \textbf{翘曲约束}：$\Delta T_{ij} = T_i - T_j \le \Delta T_{\max}$对所有相邻裸片对$(i,j)$成立——
    温差同样是$\mathbf{L}$的线性函数。
    \item \textbf{正线性单调性}：功耗越大，温度越高（物理保证，不需要近似）。
    因此我们只需要检查最坏情况（TDP + 稳态）——如果通过，所有次坏情况自动通过。
\end{itemize}

\textbf{关键观察}：无论全局功率密度（一条不等式）还是完整RC热网络+翘曲（成百上千条不等式），
功耗约束的形式始终是$\mathbf{C} \cdot \mathbf{L} \le \mathbf{d}$。
线性PDE保证：热模型的精度升级不改变约束的线性类型。
